\section{Results and Discussion}
\label{sec:results-discussion}

We first compare the networks with the linear models on Runge's function, and then with
logistic regression on MNIST. The section ends with a critical evaluation of the methods.

\subsection{Regression}

The test set has only 20 samples, and the noise variance $\sigma^2 = 0.01$ sets a floor for
the test MSE that every good model reaches. Differences in test MSE between good models are
therefore mostly noise, so we compare them by their MSE against the true function $f$ on a
dense grid, which is free of noise.

\subsubsection{Sigmoid networks and optimizers}

\Cref{fig:model-comparison} shows the main result of this part. OLS with the polynomial
degree selected in Project~1 oscillates near the ends of the interval, while both networks
follow Runge's function smoothly, using only the raw $x$ as input.

\begin{figure}[htbp]
     \centering
     \includegraphics[width=\linewidth]{figures/regression/model_comparison.pdf}
     \caption{OLS and two sigmoid networks fitted to 80 noisy samples of Runge's function.
          OLS oscillates near the ends of the interval, while the networks follow the
          function smoothly.}
     \label{fig:model-comparison}
\end{figure}

OLS has the lowest training MSE, but the networks approximate $f$ about an order
of magnitude better, as shown in \cref{tab:regression}. The polynomial spends its
flexibility on fitting the noise near the edges, where the data constrain it least.

\begin{table}[htbp]
    \centering
    \caption{Performance of OLS and the two sigmoid networks. The networks have a higher
          training MSE, but are far closer to the true function.}
    \label{tab:regression}
    {\footnotesize\setlength{\tabcolsep}{4pt}
    \begin{tabular}{@{}lccc@{}}
        \toprule
        Model & Train MSE & Test MSE & MSE vs $f$ \\
        \midrule
        OLS ($p=12$) & $0.0081$ & $0.0086$ & $0.0125$ \\
        FFNN ($1 \times 50$) & $0.0097$ & $0.0092$ & $0.0020$ \\
        FFNN ($2 \times 100$) & $0.0089$ & $0.0085$ & $0.0012$ \\
        \bottomrule
    \end{tabular}}
\end{table}

Training a sigmoid network often starts with a plateau, where the network predicts the mean of
the data and the gradient towards curved features is small~\cite{saxe_2014} (\cref{fig:training-history}).
\Cref{fig:optimizer-curves} shows that the optimizer decides how long this lasts: RMSprop and
Adam leave the plateau far sooner than plain SGD at the same learning rate. A larger learning
rate also leaves it sooner, but makes the cost noisy near the minimum, which motivates the
decaying learning rate we use for classification.

\begin{figure}[htbp]
     \centering
     \includegraphics[width=\linewidth]{figures/regression/optimizer_curves.pdf}
     \caption{Training MSE per epoch of the $1 \times 50$ sigmoid network, trained with plain
     SGD, RMSprop and Adam at batch size 16 for four learning rates. The adaptive methods
     leave the initial plateau far sooner, and larger learning rates are faster but noisier.}
     \label{fig:optimizer-curves}
\end{figure}

\subsubsection{Activation functions, depth and width}

\Cref{fig:activation-architectures} shows that small networks are best for all three
activations, and that the error grows with depth, faster for the ReLU family than for the
sigmoid. A smooth function of one variable needs little capacity, and the extra parameters of
a large network are spent on the noise.

\begin{figure*}[htbp]
     \centering
     \includegraphics[width=\linewidth]{figures/regression/activation_architectures.pdf}
     \caption{Cross-validated MSE for each activation and architecture, at the best learning
          rate $\gamma \in \{10^{-2}, 10^{-3}, 10^{-4}\}$, with Adam, 5000 epochs and a batch
          size of 32. Lighter is better. Shallow networks are best, and depth hurts the ReLU
          family most.}
     \label{fig:activation-architectures}
\end{figure*}

The best network of each activation reaches about the same test MSE
(\cref{tab:activation-comparison}), well within the noise of the small test set. Only the
MSE against $f$ separates them, in favor of the sigmoid. The ReLU networks are piecewise
linear and approximate the peak with a few straight segments, while the sigmoid gives a smooth
curve that suits Runge's function better (\cref{fig:activation-comparison}).

\begin{table}[htbp]
    \centering
    \caption{Performance of the network with the lowest cross-validated MSE for each
          activation. Only the MSE against $f$ separates them.}
    \label{tab:activation-comparison}
    {\footnotesize\setlength{\tabcolsep}{4pt}
    \begin{tabular}{@{}lccc@{}}
        \toprule
        Activation & Train MSE & Test MSE & MSE vs $f$ \\
        \midrule
        Sigmoid ($1 \times 80$) & $0.0091$ & $0.0076$ & $0.0015$ \\
        ReLU ($2 \times 10$) & $0.0083$ & $0.0080$ & $0.0022$ \\
        Leaky ReLU ($2 \times 10$) & $0.0083$ & $0.0080$ & $0.0021$ \\
        \bottomrule
    \end{tabular}}
\end{table}

The learning curves in \cref{fig:activation-curves} confirm that the large networks overfit:
their training MSE falls below the noise floor, while their validation MSE stays well above
it. The smallest networks show a much smaller gap.

\begin{figure}[htbp]
     \centering
     \includegraphics[width=\linewidth]{figures/regression/activation_curves.pdf}
     \caption{Training (dashed) and validation (solid) MSE per epoch for the smallest
          ($1 \times 10$) and largest ($8 \times 160$) network of each activation, validated
          on one fifth of the training data. The large networks fit the noise, with a
          training MSE below $\sigma^2$ and a validation MSE well above it.}
     \label{fig:activation-curves}
\end{figure}

\subsubsection{Regularization}

To counter this overfitting, we add the $\ell_1$ and $\ell_2$ penalties to the networks with
four and eight hidden layers of 40 and 80 nodes. \Cref{fig:regularization-lambda} shows the
result, and \cref{fig:regularization-heatmaps} shows the full search over $\lambda$ and $\gamma$.

\begin{figure*}[htbp]
     \centering
     \includegraphics[width=0.8\textwidth]{figures/regression/regularization_lambda.pdf}
     \caption{Training (dashed) and cross-validated (solid) MSE as functions of the penalty
          $\lambda$, at the best learning rate, with $\lambda = 0$ drawn one decade below the
          smallest $\lambda$. For the ReLU family, both penalties give a clear minimum at
          $\lambda = 10^{-3}$, while every sigmoid network collapses to the mean prediction
          already at small $\lambda$.}
     \label{fig:regularization-lambda}
\end{figure*}

For the ReLU family, both penalties give the U-shaped curve of the bias--variance trade-off.
Without a penalty, the networks fit the noise. A moderate penalty raises the training MSE
towards $\sigma^2$ and lowers the validation MSE, while a large one shrinks the weights until
the network predicts the mean. The penalty helps most for the networks that overfit most, and
brings them to the level of the best small networks without a penalty, so for the ReLU family,
a penalty can replace a careful choice of architecture. The two norms perform almost
identically, and since we use a subgradient with a finite step, $\ell_1$ does not give the
exactly sparse weights of Lasso either.

The sigmoid networks gain nothing from a penalty: their best configuration has $\lambda = 0$,
and even small penalties make them predict the mean. A possible explanation is how the
weights scale. A ReLU network can shrink the weights of one layer and grow those of another
without changing its output, and
can therefore reduce the penalty while keeping its function. A sigmoid network cannot, since
small weights put every node in the nearly linear region around $z = 0$, where the network
loses the curvature it needs for the peak. The sigmoid networks also overfit less to begin
with, so they have less to gain.

Finally, the best regularized networks approximate $f$ about five times better than the Ridge
and Lasso models (\cref{tab:regularization-comparison}).
Ridge still oscillates near $x = 1$, while Lasso underfits the peak
(\cref{fig:regularization-comparison}). The Lasso
model is weaker than it could be, since our Lasso solver did not converge for high degrees.

\begin{table}[htbp]
    \centering
    \caption{The best regularized networks and the Ridge and Lasso models of Project~1. The
          networks are about five times closer to the true function.}
    \label{tab:regularization-comparison}
    {\footnotesize\setlength{\tabcolsep}{4pt}
    \begin{tabular}{@{}lccc@{}}
        \toprule
        Model & Train MSE & Test MSE & MSE vs $f$ \\
        \midrule
        Ridge ($p = 11$, $\lambda = 10^{-6}$) & $0.0084$ & $0.0100$ & $0.0097$ \\
        ReLU ($4 \times 40$), $\ell_2$, $\lambda = 10^{-3}$ & $0.0083$ & $0.0076$ & $0.0016$ \\
        Lasso ($p = 6$, $\lambda = 10^{-2.5}$) & $0.0202$ & $0.0187$ & $0.0101$ \\
        ReLU ($4 \times 80$), $\ell_1$, $\lambda = 10^{-3}$ & $0.0086$ & $0.0085$ & $0.0022$ \\
        \bottomrule
    \end{tabular}}
\end{table}

\subsection{Image classification}

\subsubsection{Logistic regression}

Logistic regression is the network with no hidden layers, so each class score is a weighted
sum of the raw pixels. It reaches a test accuracy of about 93\%
(\cref{tab:classification-accuracy}), far above the 11\% of always guessing the most common
digit.

The errors are not spread evenly. Digits with a distinctive shape, such as 0 and 1, are
almost always recognized, while digits that share strokes, such as 3, 5 and 8, are often
confused (\cref{fig:logreg-matrix}). The weights explain why: each digit has a single
template tied to fixed pixel positions (\cref{fig:logreg-weights}), so a digit that is
shifted, slanted or written in an unusual style matches the wrong template. Handling this
variation requires features that are learned from the data, which is what we expect the
hidden layers to provide.

\begin{figure*}[htbp]
     \centering
     \includegraphics[width=\linewidth]{figures/classification/logreg_weights.pdf}
     \caption{Weights of logistic regression for each digit, reshaped to $28 \times 28$.
     Red pixels increase the score of the digit and blue pixels decrease it. Each digit is
     recognized by a single, blurred template at a fixed position.}
     \label{fig:logreg-weights}
\end{figure*}

\subsubsection{Neural networks}

\Cref{tab:classification-hyperparameter-search} summarizes the search and the selected
network. \Cref{fig:hyperparameter-search} shows that width matters more than depth: for all
activations, going from 20 to 80 nodes per layer helps, while more than four hidden layers
gives nothing for the ReLU family. The sigmoid networks behave differently. They need more
epochs, prefer large learning rates, and fail when the network is deep, which is consistent
with the vanishing gradients discussed in \cref{subsec:neural-networks}. The ReLU family
trains well over a wide range of learning rates, but fails at $\gamma = 10^{-1}$.
The best network is the widest in the grid, so a wider network might have done even better.

\begin{table}[htbp]
    \centering
    \caption{Search space and selected configuration of the classification network, chosen
    by the validation accuracy.}
    \label{tab:classification-hyperparameter-search}
    \footnotesize
    \setlength{\tabcolsep}{5pt}
    \begin{tabular}{@{}lll@{}}
        \toprule
        Hyperparameter & Search space & Selected \\
        \midrule
        Activation function & Sigmoid, ReLU, Leaky ReLU & ReLU \\
        Hidden layers & $\{4, 8, 12\}$ & $4$ \\
        Nodes per layer & $\{20, 40, 80\}$ & $80$ \\
        Learning rate $\gamma$ & $\{10^{-5}, \dots, 10^{-1}\}$ & $10^{-2}$ \\
        $\ell_2$ penalty $\lambda$ & $\{0, 10^{-4}, \dots, 10^{-1}\}$ & $10^{-4}$ \\
        \bottomrule
    \end{tabular}
\end{table}

\begin{figure*}[htbp]
     \centering
     \includegraphics[width=0.8\textwidth]{figures/classification/ffnn_architectures.pdf}
     \caption{Validation accuracy of each architecture without a penalty, at its best
     learning rate (shown in each cell). Wider networks are better for all activations,
     while depth beyond four layers does not help the ReLU family and harms the sigmoid.}
     \label{fig:hyperparameter-search}
\end{figure*}

For the selected architecture, a small penalty gives a slight improvement, while
$\lambda > 10^{-2}$ clearly hurts (\cref{fig:ffnn-regularization}). With
$60\,000$ training images and early stopping, the network does not overfit much, so there is
little for the penalty to correct. Large penalties combined with large learning rates make
training fail, and the network falls back to predicting a single digit.

\Cref{fig:ffnn-learning-curve} shows what the learning rate schedule and early stopping
add. With a constant learning rate, the validation cost stops improving early and becomes
noisy, while the decaying learning rate keeps lowering it and gives a clearly better network.
Early stopping ends training at about half the maximum number of epochs, with almost the same
accuracy as training for all of them.

\begin{figure}[htbp]
     \centering
     \includegraphics[width=\linewidth]{figures/classification/ffnn_learning_curve.pdf}
     \caption{Training (dashed) and validation (solid) cost per epoch of the selected network,
     trained for all 50 epochs with a constant and an exponentially decaying learning rate.
     The decay gives a lower validation cost, and early stopping (dotted) ends training
     where the gains have become small.}
     \label{fig:ffnn-learning-curve}
\end{figure}

\subsubsection{Comparison on the test set}

On the test set, the network reduces the error rate of logistic regression by a factor of
three (\cref{tab:classification-accuracy}). The improvement is largest for the digits
that logistic regression confused most, such as 3, 5 and 8 (\cref{fig:ffnn-matrix}). The accuracy and the
macro-averaged $F_1$ score are almost equal for all models, since MNIST is close to balanced.

Our implementations agree well with Scikit-Learn. Our logistic regression is slightly better,
and our network is about 1.5 percentage points better than \texttt{MLPClassifier} with the same
architecture, which has neither our learning rate schedule nor early stopping. This gap is
similar to the one between the constant and decaying learning rate in
\cref{fig:ffnn-learning-curve}.

\begin{table}[htbp]
    \centering
    \caption{Test accuracy and macro-averaged $F_1$ score of the selected models.}
    \label{tab:classification-accuracy}
    {\footnotesize\setlength{\tabcolsep}{4pt}
    \begin{tabular}{@{}lcc@{}}
        \toprule
        Model & Accuracy & Macro $F_1$ \\
        \midrule
        Majority class & $0.114$ & $0.020$ \\
        Logistic regression & $0.926$ & $0.925$ \\
        Logistic regression (Scikit-Learn) & $0.922$ & $0.921$ \\
        FFNN & $0.977$ & $0.977$ \\
        FFNN (Scikit-Learn) & $0.962$ & $0.961$ \\
        \bottomrule
    \end{tabular}}
\end{table}

\begin{figure}[htbp]
     \centering
     \includegraphics[width=\linewidth]{figures/classification/ffnn_confusion_matrix.pdf}
     \caption{Confusion matrix of the selected network on the validation set, normalized by
     row (\%). Compared with logistic regression (\cref{fig:logreg-matrix}), the recall is
     high for every digit.}
     \label{fig:ffnn-matrix}
\end{figure}

\subsection{Critical evaluation}\label{subsec:critical-evaluation}

For both tasks, the neural networks gave the best models, but at a much higher cost. Linear
and logistic regression have few hyperparameters, train quickly and have interpretable
weights; OLS and Ridge even have closed-form solutions, and the cost of logistic regression is
convex. A network has a non-convex cost, needs a search over the architecture, the activation,
the optimizer, the learning rate and the penalty, and each configuration must be trained for
many epochs. The networks were also more fragile: too large a learning rate or too deep a
sigmoid network made training fail, which never happened for the linear models.

For regression, a small sigmoid network approximated Runge's function best. Its advantage
over OLS and Ridge comes from the polynomial features, which oscillate near the ends of the
interval, rather than from linear regression itself; a better choice of basis would likely
narrow the gap. If cost or interpretability matters, Ridge is a reasonable choice,
but for the best fit of a function with a sharp feature, the network is better.

For classification, the network made about a third as many errors as logistic regression.
Here the features are hard to design by hand, and learning them is what makes the
difference. Logistic regression remains a strong and cheap baseline, and its weights show
directly what the model has learned.

Across both tasks, Adam and RMSprop were more robust than plain gradient descent, and ReLU
networks trained more reliably than sigmoid networks when the network was deep. A penalty
only helped where the networks overfit, which was the case for the large networks on the 80
regression samples, but hardly on the 50\,000 MNIST images. A decaying learning rate and early
stopping improved the classification network at no extra cost.

Our study also has limitations. Every configuration was trained from one set of initial
weights and evaluated on one split, so small differences between configurations are not
significant. The regression test set has only 20 samples, the best classification network is
at the edge of the search grid, and the grids are coarse. The learning rate schedule was only
used for classification. The answer to the question in \cref{sec:introduction} is therefore
qualitative: the flexibility of a network is worth its cost when good features are hard to
choose by hand, as for images, and least worth it when a simple model with suitable features
is available.
